\label{ch:strengthening}

This chapter investigates formulation-strengthening methods for the integrated aircraft routing and maintenance model. The analysis is deliberately separated from the baseline contributions of Chapters~\ref{ch:ground}--\ref{ch:hierarchy}: the corrected ground model and corrected Constraint~(13) remain the reference formulation, while each strengthening is implemented as an opt-in variant. The purpose is to determine which methods improve compactness or relaxation strength without silently changing the intended planning problem.

\section{Research questions and baseline}

The strengthening study addresses five questions:

\begin{enumerate}
    \item Can maintenance-trigger variables be removed by preprocessing without removing any feasible aircraft route?
    \item Can the C13 relaxation be tightened with interval-specific bounds without changing the integer feasible region?
    \item Do stronger maintenance-linking and maintenance-window constraints improve the linear relaxation or only increase model size?
    \item \textit{How do the available lower and upper bounds bracket the optimum, and does the strengthening tighten that bracket?}
    \item \textit{Does the combined package change exact-solver behaviour on larger instances, where the time budget actually binds?}
\end{enumerate}

The baseline is the default configuration of \texttt{MILP\_Sheduler} in \texttt{src/model.py}. It uses the corrected split form of Constraint~(13), the original maintenance-trigger domain, aggregate C8 blocking, and the global value $M=9{,}999{,}999$. All new options default to \texttt{False}; therefore, the experiments are additive and the main formulation remains reproducible.

\section{Strengthening methods}

\subsection{Aircraft-compatible maintenance domains}

The baseline creates a maintenance trigger variable $z_{ijdc}$ for every maintenance-eligible flight $i$, aircraft $j$, candidate day $d$, and check type $c$. If the assignment arc $(i,j)$ is already excluded from the sparse assignment set $X$, the corresponding maintenance variable cannot participate in a feasible solution. The sparse-domain variant therefore creates $z_{ijdc}$ only when $(i,j)\in X$. This is a preprocessing reduction, not a relaxation: an aircraft cannot perform maintenance triggered by a flight it cannot operate.

\subsection{Deferred-trigger reachability}

For a trigger flight arriving at maintenance airport $m$, a deferred check on day $d$ is possible only if the aircraft can remain at $m$ until that day. The reachability filter removes a candidate day whenever a departure from $m$ is forced after the trigger arrival and before the candidate maintenance day. Same-day candidates remain subject to the detailed C8 maintenance-window constraints. This necessary-condition test reduces the trigger domain while preserving all trigger days that can be compatible with the observed timetable.

\subsection{Strong trigger linking}

For the sparse candidate-day set $D(i,c)$, the strong-link variant adds

\begin{equation}
    \sum_{d\in D(i,c)} z_{ijdc} \leq x_{ij},
    \qquad \forall i,j,c.
    \label{eq:strength_strong_link}
\end{equation}

Equation~\eqref{eq:strength_strong_link} links every possible trigger day to the assignment variable. It is particularly relevant to deferred C/D triggers, for which the baseline link may constrain only the first candidate day. The inequality is redundant for integral assignments when the remaining model logic is sufficient, but it can eliminate unsupported fractional trigger mass from the LP relaxation.

\subsection{Interval-specific C13 bounds}

The global C13 constant is replaced, optionally, by

\begin{equation}
    M_{jdd'}^c = \max\left\{0,
    \sum_{i\in F_{d,d'}:(i,j)\in X}\tilde{t}_i
    - T_{\max}^c\right\}.
    \label{eq:strength_tight_m}
\end{equation}

The summation includes every flight compatible with aircraft $j$ in the interval, so it is an upper bound on the cumulative flight time of any feasible assignment in that interval. Replacing the global constant by equation~\eqref{eq:strength_tight_m} preserves integer-feasible solutions and can strengthen the LP relaxation. The pre-horizon C13b constraints are not modified by this experiment.

\subsection{Pairwise maintenance conflicts}

The baseline C8 formulation aggregates all blocked departures in one inequality. The pairwise variant disaggregates it into

\begin{equation}
    z_{ijdc}+x_{i'j}\leq 1,
    \qquad \forall i'\text{ blocked by the maintenance window of }(i,d,c).
    \label{eq:strength_pairwise_conflict}
\end{equation}

This formulation is at least as strong for the corresponding blocked-flight set, but may create more rows. It is therefore evaluated as a separate option rather than enabled by default.

\section{Combined formulation}

The combined Strengthening formulation enables sparse aircraft-compatible domains, reachability filtering, strong trigger linking, interval-specific C13 bounds, and pairwise maintenance conflicts. The C10 builder also skips empty station-day expressions rather than passing a trivial Boolean to Pyomo. This last change is a construction robustness correction: it does not alter a nonempty capacity constraint, but it permits long-horizon sparse experiments to distinguish model infeasibility from a model-building error.

The combined package is controlled by independent flags, so the contribution of each family can be isolated.

\section{Correctness and implementation checks}

The structural test suite verifies five properties. First, all strengthening flags are opt-in and the baseline default is unchanged. Second, every sparse maintenance variable has a corresponding assignment arc. Third, the reachability filter removes only deferred days following a forced station departure. Fourth, interval-specific $M$ values are nonnegative and finite. Fifth, pairwise conflict mode creates direct blocking rows for every blocked flight. The tests also verify that the controlled baseline and strengthened models preserve the same objective on feasible fixtures.

The complete explicitly discovered suite passes after the strengthening changes. This structural evidence is necessary but not sufficient: a smaller model is useful only if it preserves the intended solutions and improves, or at least does not materially damage, exact-solver performance.

%%==========================================================================
\section{Bounding framework}
\label{sec:strength_bounds_framework}
%%==========================================================================

\textit{The model is a minimisation problem, so its optimum is bracketed from
below and from above. The study records three bound sources separately, because
they are not interchangeable and they fail in different ways.}

\textit{\textbf{Lower bounds.} The LP relaxation value is a valid lower bound on
every integer-feasible schedule. During branch and bound the solver additionally
maintains a dual bound, which is the strongest proven lower bound at
termination. Both are recorded.}

\textit{\textbf{Upper bounds.} Any integer-feasible schedule certifies an upper
bound. Two sources are used. The solver incumbent, or primal bound, is an upper
bound whenever a feasible solution has been found. A heuristic schedule from
Chapter~\ref{ch:heuristics} is also an upper bound, but \emph{only} when it
covers every flight: a partial schedule omits flights and therefore prices an
easier problem, so its cost is recorded but explicitly not treated as a bound.}

\textit{\textbf{Residual gap.} For a dual bound $\underline{z}$ and an incumbent
$\bar z$, the reported relative gap is}
\begin{equation}
    \gamma = \frac{\bar z - \underline{z}}{|\bar z|},
    \label{eq:strength_gap}
\end{equation}
\textit{which is zero exactly when optimality has been proved. A run that
reaches its time budget without any incumbent has a lower bound but no upper
bound; that outcome is reported as such rather than scored as a large gap.}

\textit{The harness \texttt{experiments/run\_strengthening\_bounds.py} records
these quantities for every instance and both formulations. Solutions are never
auto-loaded, so a budget-limited run that finds no incumbent still reports the
dual bound it attained instead of failing.}

%%==========================================================================
\section{Benchmark design}
\label{sec:strength_bounds_design}
%%==========================================================================

\textit{Three tiers are used, because a single family cannot answer both the
bound question and the scalability question. All runs use the HiGHS backend on
the same machine with a declared per-run wall-clock budget, and the two
formulations always receive the identical instance and objective.}

\begin{description}
    \item[\textit{Small}] \textit{The 20-instance constructive A-check family of
    64 flights and 8 aircraft. This is the reference family: both formulations
    close every instance, which isolates the effect of the strengthening on
    model size and runtime.}
    \item[\textit{Medium}] \textit{Certified full-hierarchy instances of 128 and
    336 flights produced by the generator of Chapter~\ref{ch:experiments}, each
    shipped with a verified feasible certificate.}
    \item[\textit{Large}] \textit{Larger certified full-hierarchy instances of
    320 and 560 flights, where the budget is expected to bind and the residual
    gap becomes the object of study.}
\end{description}

\textit{Instances with zero-based flight identifiers were excluded. The model
indexes the cost matrix by $\text{flight}-1$, so such files silently mis-price
the first leg; only one-based, schema-validated instances enter the study.}

%%==========================================================================
\section{Model-size results}
\label{sec:strength_bounds_size}
%%==========================================================================

\textit{Table~\ref{tab:strength_bounds_model_size} reports the average model
size per tier. The reduction is driven by aircraft-compatible domain pruning and
reachability filtering; strong linking and the interval-specific $M$ values add
no variables. The saving grows with instance size, which is the expected
behaviour for a preprocessing reduction whose effect scales with the trigger
domain.}

\begin{table}[htbp]
\centering\small
\caption{Average model size per tier: baseline versus combined strengthening.}
\label{tab:strength_bounds_model_size}
\begin{tabular}{lrrrrrrr}
\toprule
Tier & Inst. & Flights & $z$ base & $z$ comb. & Cons. base & Cons. comb. & $z$ reduction \\
\midrule
Small & 20 & 64 & 2,048 & 256 & 6,598 & 3,269 & 87.5\% \\
Medium & 5 & 128 & 9,120 & 4,160 & 19,320 & 17,186 & 54.4\% \\
\bottomrule
\end{tabular}
\end{table}


\begin{figure}[htbp]
    \centering
    \includegraphics[width=0.9\textwidth]{figures/strength_bounds_zvars.png}
    \caption{\textit{Maintenance-trigger variables per instance, baseline versus combined formulation, on the largest tier.}}
    \label{fig:strength_bounds_zvars}
\end{figure}

%%==========================================================================
\section{Bound results}
\label{sec:strength_bounds_results}
%%==========================================================================

\textit{Table~\ref{tab:strength_bounds_summary} summarises the exact-solve
outcome per tier and formulation.}

\begin{table}[htbp]
\centering\small
\caption{Exact-solve outcome per tier: proven optimality, residual MIP gap and runtime.}
\label{tab:strength_bounds_summary}
\begin{tabular}{llrrrrrr}
\toprule
Tier & Variant & Inst. & Optimal & Mean gap (\%) & Max gap (\%) & Mean time (s) & Median time (s) \\
\midrule
Small & baseline & 20 & 20/20 & 0.000 & 0.000 & 1.69 & 1.69 \\
Small & combined & 20 & 20/20 & 0.000 & 0.000 & 1.00 & 1.00 \\
Medium & baseline & 5 & 5/5 & 0.000 & 0.000 & 4.92 & 5.03 \\
Medium & combined & 5 & 5/5 & 0.000 & 0.000 & 3.21 & 3.22 \\
\bottomrule
\end{tabular}
\end{table}


\textit{The first observation concerns the reference tier and is negative for
the bound question. On that family the LP relaxation is already integral: the LP
value, the dual bound and the incumbent coincide on every instance, so the
integrality gap is zero and there is nothing for a cut to close. On the
full-hierarchy tier the picture is different, and a strictly positive
integrality gap does appear on part of the family, as
Table~\ref{tab:strength_bounds_bracket} shows. Whether the relaxation is tight
is therefore a property of the instance family, not of the formulation.}

\textit{The decisive observation is that the baseline and combined LP bounds
agree on every instance of every tier tested here, to solver tolerance. The
combined package removes variables and rows that are unsupported at the LP
optimum, but it does not move that optimum. Its measurable benefit is therefore
compactness and solve time, and the strengthening should not be presented as a
bound-tightening device on these families.}

\textit{This is a correction to an earlier reading of the same study. An LP bound
improvement reported previously was obtained under a different solver and check
configuration, before the maintenance-window and horizon corrections of
Chapters~\ref{ch:maintenance}--\ref{ch:hierarchy}. It is not reproducible under
the current formulation and is therefore withdrawn.}

\textit{Table~\ref{tab:strength_bounds_bracket} reports the per-instance bracket
and Figure~\ref{fig:strength_bounds_bracket} plots it. The shaded band runs from
the LP lower bound to the incumbent, so its width is the integrality gap; the
dual bound at termination shows how much of that band the search actually
closed.}

\begin{table}[htbp]
\centering\small
\caption{Bound bracket under the combined formulation. The integrality gap separates the LP bound from the incumbent; the residual gap is what remained unproved at termination. Heuristic coverage below 100\% means the heuristic cost is not a valid upper bound.}
\label{tab:strength_bounds_bracket}
\begin{tabular}{lrrrrrr}
\toprule
Instance & LP bound & Dual bound & Incumbent & Integrality gap (\%) & Residual gap (\%) & Heur. coverage (\%) \\
\midrule
DataCplex\_density=0.65\_p=8\_h=14\_test\_0 & 29,738.0 & 29,738.0 & 29,738.0 & 0.000 & 0.000 & 84.4 \\
DataCplex\_density=0.65\_p=8\_h=14\_test\_1 & 29,736.0 & 29,736.0 & 29,736.0 & 0.000 & 0.000 & 84.4 \\
DataCplex\_density=0.65\_p=8\_h=14\_test\_10 & 29,774.0 & 29,774.0 & 29,774.0 & 0.000 & 0.000 & 84.4 \\
DataCplex\_density=0.65\_p=8\_h=14\_test\_11 & 29,774.0 & 29,774.0 & 29,774.0 & 0.000 & 0.000 & 84.4 \\
DataCplex\_density=0.65\_p=8\_h=14\_test\_12 & 29,799.0 & 29,799.0 & 29,799.0 & 0.000 & 0.000 & 84.4 \\
DataCplex\_density=0.65\_p=8\_h=14\_test\_13 & 29,769.0 & 29,769.0 & 29,769.0 & 0.000 & 0.000 & 84.4 \\
DataCplex\_density=0.65\_p=8\_h=14\_test\_14 & 29,746.0 & 29,746.0 & 29,746.0 & 0.000 & 0.000 & 84.4 \\
DataCplex\_density=0.65\_p=8\_h=14\_test\_15 & 29,767.0 & 29,767.0 & 29,767.0 & 0.000 & 0.000 & 84.4 \\
DataCplex\_density=0.65\_p=8\_h=14\_test\_16 & 29,750.0 & 29,750.0 & 29,750.0 & 0.000 & 0.000 & 84.4 \\
DataCplex\_density=0.65\_p=8\_h=14\_test\_17 & 29,716.0 & 29,716.0 & 29,716.0 & 0.000 & 0.000 & 84.4 \\
DataCplex\_density=0.65\_p=8\_h=14\_test\_18 & 29,729.0 & 29,729.0 & 29,729.0 & 0.000 & 0.000 & 84.4 \\
DataCplex\_density=0.65\_p=8\_h=14\_test\_19 & 29,763.0 & 29,763.0 & 29,763.0 & 0.000 & 0.000 & 84.4 \\
DataCplex\_density=0.65\_p=8\_h=14\_test\_2 & 29,776.0 & 29,776.0 & 29,776.0 & 0.000 & 0.000 & 84.4 \\
DataCplex\_density=0.65\_p=8\_h=14\_test\_3 & 29,744.0 & 29,744.0 & 29,744.0 & 0.000 & 0.000 & 84.4 \\
DataCplex\_density=0.65\_p=8\_h=14\_test\_4 & 29,777.0 & 29,777.0 & 29,777.0 & 0.000 & 0.000 & 84.4 \\
DataCplex\_density=0.65\_p=8\_h=14\_test\_5 & 29,730.0 & 29,730.0 & 29,730.0 & 0.000 & 0.000 & 84.4 \\
DataCplex\_density=0.65\_p=8\_h=14\_test\_6 & 29,797.0 & 29,797.0 & 29,797.0 & 0.000 & 0.000 & 84.4 \\
DataCplex\_density=0.65\_p=8\_h=14\_test\_7 & 29,759.0 & 29,759.0 & 29,759.0 & 0.000 & 0.000 & 84.4 \\
DataCplex\_density=0.65\_p=8\_h=14\_test\_8 & 29,771.0 & 29,771.0 & 29,771.0 & 0.000 & 0.000 & 84.4 \\
DataCplex\_density=0.65\_p=8\_h=14\_test\_9 & 29,784.0 & 29,784.0 & 29,784.0 & 0.000 & 0.000 & 84.4 \\
sma1\_p8\_h8\_r1 & 161,149.6 & 161,149.6 & 161,149.6 & 0.000 & 0.000 & 34.4 \\
sma1\_p8\_h8\_r2 & 161,987.1 & 161,987.1 & 161,987.1 & 0.000 & 0.000 & 31.2 \\
sma1\_p8\_h8\_r3 & 159,245.3 & 159,245.3 & 159,245.3 & -0.000 & 0.000 & 39.1 \\
sma1\_p8\_h8\_r4 & 160,758.8 & 161,360.8 & 161,360.8 & 0.373 & 0.000 & 38.3 \\
sma1\_p8\_h8\_r5 & 161,586.0 & 161,904.1 & 161,904.1 & 0.197 & 0.000 & 39.8 \\
\bottomrule
\end{tabular}
\end{table}


\begin{figure}[htbp]
    \centering
    \includegraphics[width=0.95\textwidth]{figures/strength_bounds_bracket.png}
    \caption{\textit{Bound bracket under the combined formulation. The band spans the LP lower bound to the incumbent, and the dashed series is the dual bound proved at termination.}}
    \label{fig:strength_bounds_bracket}
\end{figure}

\textit{The heuristic upper bound is reported only where coverage is complete,
and on these families it never is. The greedy-plus-insertion heuristic leaves a
large part of the timetable unassigned, and the coverage column of
Table~\ref{tab:strength_bounds_bracket} records how much. Its cost is therefore
not a valid upper bound and is excluded from the bracket. The point is worth
stating plainly because the omission is easy to misread: a partial schedule can
look far cheaper than the optimum precisely because it declines to fly the legs
it cannot accommodate, so treating that cost as a bound would be a modelling
error rather than a conservative approximation.}

\begin{figure}[htbp]
    \centering
    \includegraphics[width=0.95\textwidth]{figures/strength_bounds_gap.png}
    \caption{\textit{Residual MIP gap per instance on the largest tier. Bars at zero indicate optimality proved within the budget.}}
    \label{fig:strength_bounds_gap}
\end{figure}

%%==========================================================================
\section{Larger-model statistics}
\label{sec:strength_bounds_large}
%%==========================================================================

\textit{On the larger tier the model dimensions diverge sharply between the two
formulations and the time budget begins to bind. Where the budget binds, the
comparison is no longer about speed but about what each formulation can prove
within a fixed effort: how many instances are closed, what dual bound is
attained, and whether any incumbent is found at all.}

\textit{Table~\ref{tab:strength_bounds_scaling} reports model size on larger
certified instances of 320 and 560 flights. Two effects are visible and they
pull in opposite directions. The maintenance-trigger domain shrinks on every
instance, and the saving grows with the horizon. The constraint count, however,
does not fall uniformly: on the 560-flight, eight-aircraft instances it drops
substantially, whereas on the 320-flight instances with a twenty-aircraft fleet
it rises, because disaggregating the maintenance-window conflicts creates one
row per blocked flight and that expansion scales with the fleet. The combined
package is therefore not uniformly smaller; it trades a much smaller integer
domain against a possibly larger constraint matrix, and which effect dominates
depends on fleet size.}

\begin{table}[htbp]
\centering\small
\caption{Model-size scaling on the larger certified instances. The trigger domain shrinks throughout, but the constraint count can rise when pairwise maintenance conflicts are disaggregated over a large fleet.}
\label{tab:strength_bounds_scaling}
\begin{tabular}{lrrrrrrrr}
\toprule
Instance & Flights & Aircraft & $z$ base & $z$ comb. & $z$ drop (\%) & Cons. base & Cons. comb. & Cons. change (\%) \\
\midrule
med1\_p20\_h8\_r1 & 320 & 20 & 59,200 & 24,800 & 58.1 & 155,624 & 191,584 & $+$23.1 \\
med1\_p20\_h8\_r2 & 320 & 20 & 59,200 & 24,800 & 58.1 & 155,624 & 195,304 & $+$25.5 \\
med1\_p20\_h8\_r3 & 320 & 20 & 59,200 & 24,800 & 58.1 & 155,624 & 191,064 & $+$22.8 \\
sma3\_p8\_h35\_r1 & 560 & 8 & 127,552 & 17,536 & 86.3 & 214,503 & 86,711 & -59.6 \\
sma3\_p8\_h35\_r2 & 560 & 8 & 127,552 & 17,536 & 86.3 & 214,503 & 86,599 & -59.6 \\
sma3\_p8\_h35\_r3 & 560 & 8 & 127,552 & 17,536 & 86.3 & 214,503 & 86,407 & -59.7 \\
\bottomrule
\end{tabular}
\end{table}


\textit{This also explains why the largest configurations were excluded from the
solve study. At 560 flights the baseline model carries over 127{,}000
maintenance-trigger variables, and neither formulation returned an incumbent
within a short budget; only the combined model returned a dual bound at all.
Reporting a residual gap in that regime would be meaningless, since one side of
the bracket does not exist, so those runs are reported as model-size evidence
rather than as bound evidence.}

\textit{Figure~\ref{fig:strength_bounds_runtime} reports the paired exact
runtimes on a logarithmic scale. Runtimes are wall-clock and include transferring
the model to the solver, which is itself substantial at this size, so the smaller
combined model is favoured twice: once in transfer and once in search.}

\begin{figure}[htbp]
    \centering
    \includegraphics[width=0.95\textwidth]{figures/strength_bounds_runtime.png}
    \caption{\textit{Paired exact-solve wall-clock runtimes on the largest tier, logarithmic scale.}}
    \label{fig:strength_bounds_runtime}
\end{figure}

%%==========================================================================
\section{Interpretation and limitations}
%%==========================================================================

\textit{The evidence supports a narrower claim than uniform improvement.}

\textit{The compactness result is robust and grows with instance size. The
combined package removes the large majority of maintenance-trigger variables
without removing any feasible aircraft route, and that property is verified
structurally by the test suite rather than inferred from runtime.}

\textit{The bound result is conditional. On the small reference family the LP
relaxation is integral, so no strengthening can tighten the bound. On the larger
instances the relevant question changes: within a fixed budget the smaller model
is the one more likely to return a usable dual bound, so the reported bracket is
evidence about bounds attainable under a budget, not proof of a structurally
tighter relaxation.}

\textit{The upper-bound side remains the weaker half of the bracket. The solver
incumbent is the only reliable upper bound in these experiments, because the
current heuristics do not achieve complete coverage on the constructive family.
Producing a complete feasible schedule cheaply would improve the bracket more
than further reductions of the trigger domain.}

\textit{Finally, the scope is limited to the declared tiers, the HiGHS backend
and the stated budgets. Instances that are infeasible under both formulations
are excluded from efficacy claims and retained only as diagnostic cases, and no
statistical generalisation beyond these families is asserted.}

\section{Reproducibility}

\textit{Every number in this chapter is generated. The harness
\texttt{experiments/run\_strengthening\_bounds.py} writes the raw per-run CSV
files, and \texttt{experiments/analyze\_strengthening\_bounds.py} derives the
summary CSV, the LaTeX tables included above and the figures. No value in this
chapter is transcribed by hand, so re-running the two scripts regenerates the
tables and figures together, and any change in the underlying runs is reflected
automatically.}
